\section{Institutional setting and closure samples}
\label{sec:setting}

The question this paper answers is whether losing a particular inpatient
provider changed acute mortality in the places that depended on it. Both
the question and the measurement problem inside it come from three
features of Brazilian health care. Patients travel across municipal lines
for hospital care on a scale that makes geographic proximity a weak guide
to who relies on a given hospital. The psychiatric reform that dominates
the closure record shut specialized hospitals on a federally driven
schedule, and that schedule is what makes closure timing usable for
identification. The closures themselves, read off the federal
establishment register, define the sample. I take the three in turn.

% =====================================================================
% §2.1 — SUS regionalization and cross-border patient flow
% =====================================================================
\subsection{SUS regionalization and cross-border patient flow}
\label{sec:setting-sus}

Brazil's Sistema Único de Saúde (SUS), created by the $1988$ Constitution
and built out by Laws~$8.080$ and $8.142$ of $1990$, is a tax-financed
system meant to cover the entire resident population. About $71\%$ of
Brazilians use it as their only source of care, and another $25\%$ use it
alongside private insurance \citep{macinkoHarris2015}. Inpatient care
runs through a mix of municipal, state, philanthropic, and contracted
private hospitals, arranged in a nested geography. Each of Brazil's
$5{,}570$ municipalities is a management unit; municipalities group into
$438$ Regiões de Saúde (CIRs) meant to coordinate flows across levels of
care; and the $27$ Federation Units operate or contract the referral
hospitals at the top.

\paragraph{Flow does not follow geography.} The architecture prescribes a
regional hierarchy. It does not enforce one. Patients cross municipal and
CIR lines as a matter of routine, most of all for specialized care. Of
the $\valNadmissions$~million SUS admissions in the pooled patient-flow
graph ($2010$--$2024$), $\valCrossBorderPct\%$ are recorded at a hospital
outside the patient's home municipality on a volume-weighted basis;
weighting municipalities equally, the mean outflow share is $70\%$, which
says that small municipalities lean on care elsewhere far more than the
national volume figure lets on. In Sergipe, Pernambuco, and the Distrito
Federal the cross-border rate tops $55\%$. These trips are not noise
around a local-care norm. They are how much of the SUS population reaches
a hospital bed at all, and they are what breaks the distance proxy: a
patient whose usual referral hospital sits one CIR over is coded
unexposed to its closure under a nearest-hospital rule, even though that
hospital is where she actually went. The distance proxy standard in the
closure literature \citep{avdic2016, carroll2019, joynt2015,
lindrooth2018} assumes flow tracks road geography. That assumption holds
up where local capacity is thick, as in Sweden or much of the NHS, and is
shakier even in U.S. Medicare, where \citet{wang2020hrr} find patient
flows drifting away from the Hospital Referral Region lines the Dartmouth
Atlas has drawn since the early $1990$s. Brazil's $\valCrossBorderPct\%$
runs larger than any of these, and it is exactly where a flow-based
correction should bite hardest.

% =====================================================================
% §2.2 — Reforma Psiquiátrica as federal closure shock
% =====================================================================
\subsection{The Reforma Psiquiátrica and the source of closure timing}
\label{sec:setting-psiquiatrica}

The closure record is dominated by a long-running overhaul of
mental-health care, and that overhaul is both where the identifying
variation comes from and why the mortality question is politically
loaded. Lei~$10.216$ of $2001$ mandates deinstitutionalization:
asylum-style psychiatric hospitals (\emph{hospitais psiquiátricos}) are
to give way to community outpatient networks (Centros de Atenção
Psicossocial, CAPS) and supported residences (Serviços Residenciais
Terapêuticos). \citet{diasFontes2024} study the CAPS rollout and find it
expanded outpatient care and cut psychiatric hospitalizations without
lowering mental-health mortality. This paper takes the other side of the
same reform, the closures, and asks what happened in the catchments that
had relied on those inpatient beds. Two federal levers drove the
closures, and they reinforced each other.

\paragraph{PNASH inspection cycles.} The Programa Nacional de Avaliação
dos Serviços Hospitalares (PNASH) inspected every accredited psychiatric
hospital in waves---$2002$--$2003$, $2007$--$2008$, $2011$--$2012$, and
$2015$--$2016$---scoring infrastructure, staffing, and care processes. A
facility scoring below roughly $61\%$ on the composite index drew an
indication for de-accreditation under Portaria SAS/MS~$501$ of $2007$.
The schedule is the reason to expect closure timing to be loosely tied,
at most, to municipality-specific mortality shocks: the inspection
calendar is set in Brasília, not chosen locally in response to a town's
demand or supply conditions. The primary sample anchors to that calendar
(closures within $\pm 1$ year of a cycle;
Section~\ref{sec:method-identification}).

\paragraph{Reimbursement erosion.} The SUS daily rate for psychiatric
inpatient care sat near R\$$30$ per patient-day from the early $2000$s
and was never indexed to inflation; across the panel window its real
value fell by about half. The squeeze pushed facilities into insolvency
at scale, again on a clock set by federal reimbursement policy rather
than by any single municipality's finances or health needs.

Both levers apply the same federal thresholds and the same reimbursement
schedule to all $5{,}570$ municipalities, and both produce decisions at
the level of the institution, not the municipality. None of this makes
closure timing random. It does make a municipality's own health shocks a
less natural explanation for when a psychiatric hospital closed than the
PNASH cycle and the reimbursement squeeze. The design therefore treats
timing as plausibly orthogonal to local shocks, conditional on the
PNASH-anchored sample, the pre-trend checks, and the robustness exercises
that follow. Cross-referencing the closures against the homologated PNASH
results bears out the mixed-driver picture: of the $\valNclosuresPnash$
psychiatric closures, $\valPnashDescred$ were formally indicated for
de-accreditation in the $2012$--$2014$ cycle (Portaria SAS/MS~$1.727$/$2016$),
while $\valPnashClassified$ had been classified as adequate and closed
anyway---closure here is driven by reimbursement pressure alongside, not
only by, quality de-accreditation. I also classified each event's declared
closure motive from documentary sources manually. That classification organizes robustness
sub-samples and is reported in the replication notes, but its external
validation is incomplete, so it does no identifying work.

% =====================================================================
% §2.3 — Closure events 2010–2024
% =====================================================================
\subsection{\texorpdfstring{Closure events, $2010$--$2024$}{Closure events, 2010--2024}}
\label{sec:setting-closures}

Closures come from the Cadastro Nacional de Estabelecimentos de Saúde
(CNES), the federal establishment register. I flag every establishment
whose reporting stops before the window ends. The raw list is mostly
registry debris---small clinics, day units, code reassignments---so a
cumulative filter (F1--F5; online appendix) pares it to
$\valNexoClosures$ closures that clear five conditions: closure inside
the non-pandemic windows $2012$--$2018$ or $2022$--$2023$; a
hospital-grade establishment type (\texttt{tp\_unid} $\in
\{05, 07, 15, 62\}$); at least $30$ SUS-contracted beds and at least
$100$ admissions in the year before closure; and no admission drop above
$50\%$ in that final year, which screens out hospitals that had already
emptied out under falling demand. Of the $\valNexoClosures$,
$\valNspecClosures$ are specialized hospitals (\texttt{tp\_unid} $= 07$,
mostly psychiatric), $\valNgenClosures$ are general (\texttt{tp\_unid} $=
05$), and one is a day unit. The Southeast accounts for $62\%$ of the
events and the Northeast for $25\%$; the busiest year is $2012$, with
$11$ events, an echo of the early audit cycles.

\paragraph{The causal sample is narrower than the measurement sample.}
The mortality analysis does not use all $\valNexoClosures$ closures. It
uses the PNASH-anchored psychiatric subset, where the timing argument is
strongest and where the suicide and self-harm margins line up with the
nature of the shock. The primary causal sample is the
$\valNclosuresPnash$ psychiatric closures falling within $\pm 1$ year of
a PNASH cycle; the $\valNspecClosures$ statistically filtered specialized
closures and a wider $\valNclosuresPsymax$-closure psychiatric set serve
as robustness checks (Section~\ref{sec:robustness}). The hand-coded
motive labels enter only as additional robustness sub-samples and are
cross-tabulated against facility type in the replication notes. They
document the events; they do not certify any one of them as causal.

\subsection{Four samples, four jobs}
\label{sec:setting-samples}

The paper leans on four samples, each for a different claim, and keeping
them apart matters for reading the results. The F5 sample of
$\valNexoClosures$ economically meaningful closures is the measurement
sample: it carries the flow-versus-distance comparison and the general
first-stage diagnostics. The $\valNclosuresPnash$ PNASH-anchored
psychiatric closures are the causal sample for suicide and self-harm. The
$\valNgenClosures$ general-hospital closures and the lone hospital-day
closure are contrast cases, used only to ask whether flow-defined
catchments matter outside psychiatric reform. Closures that pass the bed
and volume screens but fail F5 on a large pre-closure decline are
diagnostic cases, held out of the main design because for them the
official closure date is a poor stand-in for an abrupt loss of inpatient care.
Table~\ref{tab:sample-architecture-compact} sets out the four and the
role each one plays.

\begin{table}[!htbp]\centering
\caption{Closure samples and their role in the paper}
\label{tab:sample-architecture-compact}
\small
\resizebox{\textwidth}{!}{%
\begin{tabular}{llrrl}
\toprule
Sample & Definition & Closures & Flow-exposed munis & Role \\
\midrule
F5 measurement sample & Economically meaningful hospital closures passing F1--F5 & $\valNexoClosures$ & 130 & Measurement: flow versus distance \\
PNASH psychiatric causal sample & Psychiatric closures within \(\pm 1\) year of PNASH cycles & $\valNclosuresPnash$ & $\valNtrPnash$ & Causal mortality analysis \\
General-hospital contrast & F5 closures with general-hospital type & 18 & 24 & Contrast evidence \\
Hospital-day contrast & F5 closures with hospital-day type & 1 & 1 & Contrast evidence \\
Failed-F5 diagnostic sample & Pass F4 but fail F5 due to pre-closure volume decline & 30 & 26 & Endogeneity/filter diagnostic \\
\bottomrule
\multicolumn{5}{p{0.95\textwidth}}{\footnotesize Notes: F5 is the measurement-sample filter for economically meaningful hospital closures. The PNASH psychiatric causal sample prioritizes institutional anchoring to federal inspection cycles; stricter and broader psychiatric samples are reported as robustness checks.}\\
\end{tabular}
}%
\end{table}

